% Capítulo 7 - Conclusão

\chapter{CONCLUSÃO}
\label{cap:conclusao}

Com base na implementação e nos testes realizados, conclui-se que é tecnicamente viável construir um servidor \textit{NAS} doméstico de baixo custo baseado em \textit{Radxa ROCK 4B+}, \textit{OpenMediaVault}, \textit{Nextcloud} e \textit{Tailscale} para cenários de armazenamento, compartilhamento, sincronização e acesso remoto de uso doméstico.

A arquitetura final com dois \textit{SSDs} mostrou-se funcional e administrável. O \textit{SMB} forneceu acesso local simples e desempenho adequado ao uso cotidiano; o \textit{Nextcloud} acrescentou recursos de nuvem privada sem impedir taxas úteis de transferência; e o \textit{Tailscale} permitiu acesso remoto privado por rede móvel. A utilização do \textit{Tailscale} no \textit{host} com \texttt{netfilter-mode=off} foi necessária para contornar incompatibilidades observadas no ambiente e passou a compor a configuração final validada.

Com o cliente cabeado em \textit{Gigabit}, o \textit{NAS} atendeu aos valores de referência de leitura e escrita definidos para o trabalho, com cerca de 81 \textit{MB/s} e 76 \textit{MB/s} no \textit{SMB}. Nos testes consolidados, o desempenho ficou abaixo do potencial teórico de uma rede \textit{Gigabit}, especialmente porque o cliente permaneceu em \textit{Wi-Fi}; ainda assim, o cenário com \textit{NAS} em \textit{Ethernet} apresentou melhora expressiva, indicando que o cabeamento do servidor é a configuração recomendada. Para o contexto doméstico avaliado, as taxas obtidas foram suficientes para documentos, fotos, mídia, \textit{backups} pessoais e sincronização.

A solução também demonstrou viabilidade administrativa: após a implantação inicial, o uso diário pode ser realizado prioritariamente pelas interfaces \textit{web} do \textit{OpenMediaVault} e do \textit{Nextcloud}, enquanto o terminal fica restrito a implantação, diagnóstico e manutenção. A criação de \textit{backups} automáticos de configuração reduziu o risco operacional e melhorou a recuperabilidade do ambiente.

Na comparação econômica, o custo-base efetivamente pago do protótipo, de R\$ 625,48 e excluindo os \textit{SSDs}, permaneceu cerca de 73\% inferior ao de \textit{NAS} comerciais de duas baias também vendidos sem discos. Essa comparação é particularmente adequada à arquitetura final do estudo, que utiliza dois \textit{SSDs}, assim como os modelos comerciais analisados oferecem duas baias. A solução, entretanto, exige maior conhecimento técnico e não oferece o mesmo nível de integração física e suporte comercial. A comparação com \textit{NAS} comerciais confirma que a proposta não busca equivalência de desempenho sequencial máximo: modelos como \textit{Synology DS223j} e \textit{QNAP TS-233} podem se aproximar do limite de uma interface \textit{Gigabit} em condições controladas, e o protótipo apresentou taxas menores nos cenários com \textit{Wi-Fi}, porém adequadas ao uso doméstico avaliado.

Do ponto de vista de segurança e disponibilidade, a solução não deve ser interpretada como substituta universal da nuvem pública. Seu diferencial está no controle físico dos dados, no menor grau de dependência de terceiros e no acesso remoto privado via \textit{Tailscale}. Em troca, o usuário assume responsabilidade direta por manutenção, atualização e \textit{backup}. Para o perfil doméstico analisado, esse equilíbrio entre autonomia, custo e desempenho sustenta a conclusão de viabilidade técnica.

Como trabalhos futuros, sugerem-se a medição elétrica direta do conjunto, a repetição do teste de estresse de 24 horas, a repetição completa dos testes locais com cliente e servidor cabeados, a avaliação de uso simultâneo por vários clientes e a evolução da automação da instalação e da recuperação. Esses itens ampliariam as evidências, mas não são necessários para sustentar a conclusão principal, pois a viabilidade foi demonstrada pelo protótipo real e pelas evidências experimentais coletadas.
